% =====================================================================
%  CHAPITRE 5 - RECUL SUR LE STAGE
% =====================================================================
\chapter{Recul sur le stage}
\label{chap:recul}

Après avoir détaillé les travaux réalisés, ce chapitre prend du recul sur ce que
cette expérience m'a réellement apporté. Il ne revient pas sur les solutions
techniques mises en \oe uvre, mais sur la manière dont ma façon de travailler a
évolué, sur les compétences que j'ai construites et sur celles qu'il me reste à
approfondir.

\section{Apports techniques}
\label{sec:apports-techniques}

L'apport le plus déterminant du stage tient moins aux technologies employées
qu'à la nature du terrain sur lequel j'ai travaillé. Jusqu'ici, mes projets
étaient conçus depuis une page blanche~: je définissais moi-même le modèle de
données, l'architecture et les conventions, et le résultat était évalué à sa
livraison. \csbag m'a placé dans la situation inverse. J'ai hérité d'une
application en service, avec des données réelles saisies pendant plusieurs
années, des choix antérieurs que je n'avais pas faits, et des utilisateurs qui
continuaient à s'en servir pendant que je la modifiais. La question n'est alors
jamais seulement « quelle est la bonne conception~? », mais « comment atteindre
la bonne conception sans casser ce qui fonctionne~? ». La refonte du modèle de
données l'a illustré avec le plus de force~: je connaissais les formes normales
en théorie, j'ignorais ce que représente concrètement leur application sur des
données déjà dégradées. Lors de la migration~V019, les doublons présents en base
empêchaient la création des clés primaires, et il a fallu nettoyer avant de
contraindre, situation qu'un exercice académique, où les données sont propres
par construction, ne permet pas de rencontrer. J'en ai retiré une conviction
durable~: la qualité d'un modèle de données ne se juge pas sur son schéma, mais
sur ce qu'il rend possible ou impossible aux utilisateurs. Un champ de texte
libre ne produit pas immédiatement de bug~; il produit, au bout de quelques
années, un inventaire sur lequel plus aucun filtre n'est fiable. C'est également
l'analyse en amont qui m'a le plus fait progresser~: prendre le temps de
requêter les données réelles avant de coder m'a permis de formuler un problème
que personne n'avait exprimé, et de le documenter suffisamment pour qu'il soit
accepté.

La migration du socle technique m'a appris une autre forme de rigueur. Le
réflexe naturel aurait été de sauter directement d'Angular~14 à Angular~20~; le
choix d'une progression palier par palier, plus long, m'a montré la valeur d'une
méthode qui isole les causes de régression au lieu de les cumuler. Plus
largement, j'ai découvert que la maintenance d'un socle n'est pas une tâche
secondaire mais une condition de la sécurité de l'application, puisqu'une
version sortie du support ne reçoit plus aucun correctif.

Les interventions sur Elium Sauvegarde et Access Report TTC~GO ont
considérablement fait évoluer ma manière de diagnostiquer un problème. Dans les
deux cas, l'application ne produisait aucune erreur~: Elium échouait
globalement à cause d'une anomalie sur un seul article, et Access Report livrait
des rapports incomplets sans qu'aucune trace ne soit journalisée. J'ai appris à
me méfier de l'absence de message d'erreur, à raisonner par élimination plutôt
qu'à modifier le code au hasard, et surtout à ne pas m'arrêter au premier
symptôme corrigé. Sur Access Report, la distinction entre le contournement mis
en place pour débloquer rapidement l'exploitation et la correction de fond
apportée ensuite à la bibliothèque partagée m'a fait comprendre qu'un correctif
a une portée~: résoudre le problème pour soi ou le résoudre pour tous ceux qui
utilisent le même composant ne demandent ni le même effort, ni le même niveau de
précaution. C'est aussi sur ces missions que j'ai mesuré l'importance de
l'observabilité~: un traitement qui ne dit rien de ce qu'il a fait est un
traitement que personne ne peut valider.

L'exigence de qualité imposée par la chaîne d'intégration continue a été, au
début, difficile à accepter. Voir une livraison bloquée pour une complexité
cognitive excessive ou pour une description manquante sur un tableau HTML m'a
d'abord semblé formaliste. J'ai progressivement compris que ces règles rendent
explicite une exigence que je n'appliquais jusqu'ici qu'au jugé, et qu'un code
qui passe ces contrôles est un code qu'un autre développeur pourra reprendre.
De la même façon, la découverte de la gestion des secrets et des mécanismes
d'autorisation m'a fait passer d'une conception défensive de la sécurité, où
l'on protège l'interface, à une conception plus juste, où l'on considère que
l'interface peut être contournée et que le contrôle doit exister côté serveur.

Sur le plan de l'autonomie, l'évolution a été nette. Les premières semaines,
je validais chaque orientation avant de l'engager. À partir du deuxième lot,
j'arrivais aux points techniques avec une proposition argumentée et ses
alternatives, et la discussion portait sur le choix plutôt que sur la direction
à prendre. Cette autonomie reste toutefois inégale selon les domaines~: je suis
aujourd'hui à l'aise sur la conception applicative et la donnée, nettement
moins sur les aspects d'exploitation, où j'ai surtout observé et exécuté des
procédures existantes.

\section{Apports humains et organisationnels}
\label{sec:apports-humains}

Avant ce stage, ma représentation du travail en équipe reposait sur les projets
tuteurés~: un groupe d'étudiants au même niveau, un objectif commun, et une
répartition des tâches décidée entre pairs. L'équipe que j'ai rejointe
fonctionne autrement. Chacun y occupe un rôle défini, avec des attentes et une
temporalité propres, et une part importante du travail consiste à s'articuler
avec les autres plutôt qu'à avancer seul. J'ai dû apprendre à identifier le bon
interlocuteur selon la question~: le chef de projet pour un arbitrage de
périmètre, le référent technique pour une orientation de conception, l'analyste
fonctionnelle pour lever une ambiguïté dans une spécification.

La relation avec l'analyste fonctionnelle a été particulièrement formatrice.
Recevoir un cahier de retours documenté après chaque mise en recette m'a
d'abord été inconfortable~: un écran que je considérais comme terminé revenait
avec une liste d'anomalies. J'ai compris assez vite que ces retours ne portaient
pas un jugement sur mon travail, mais sur l'écart entre ce que j'avais compris
du besoin et ce qui était attendu. Cette prise de conscience a modifié ma façon
de travailler en amont~: j'ai pris l'habitude de reformuler les points ambigus
dès la réunion de cadrage plutôt que de les interpréter, ce qui a réduit le
nombre d'allers-retours au fil des lots.

Le moment le plus structurant du stage a été celui où j'ai présenté à l'équipe
mon diagnostic de l'existant et proposé d'élargir le périmètre de la mission.
Il ne suffisait pas d'avoir raison techniquement~: il fallait démontrer que ces
chantiers, absents des spécifications, conditionnaient la valeur des
fonctionnalités demandées, et que le coût supplémentaire était justifié. J'ai
appris à formuler un problème technique dans des termes compréhensibles par des
interlocuteurs qui n'allaient pas lire le code, en m'appuyant sur des exemples
concrets tirés des données réelles plutôt que sur des principes de conception.
C'est probablement la compétence la plus transférable que j'ai développée.

Les rituels agiles ont eu un effet que je n'avais pas anticipé. Le point
quotidien oblige à rendre son travail intelligible en quelques phrases, et à
signaler un blocage au lieu de s'obstiner~: j'ai mis plusieurs semaines à
accepter que remonter une difficulté rapidement n'est pas un aveu de faiblesse,
mais le comportement attendu. Le point technique hebdomadaire, lui, m'a habitué
à défendre un choix devant quelqu'un de plus expérimenté, et à distinguer une
critique de ma solution d'une critique de mon travail.

Le fait de mener trois applications en parallèle, sur des piles techniques et
des problématiques différentes, m'a confronté à la gestion des priorités.
Basculer d'une évolution fonctionnelle planifiée vers un incident remonté par un
utilisateur suppose d'accepter de suspendre une tâche en cours, ce que je
faisais mal au départ~: j'avais tendance à vouloir terminer avant de changer de
sujet, y compris lorsque l'urgence justifiait l'inverse. Le suivi formalisé des
tâches et les échanges avec le chef de projet m'ont aidé à hiérarchiser selon
l'impact plutôt que selon mon confort.

Enfin, la participation aux réunions concernant d'autres projets du service, aux
mises en production de renouvellement de secrets et au traitement des problèmes
remontés par les utilisateurs m'a donné une vision que le seul développement ne
m'aurait pas apportée. J'ai découvert qu'une application vit bien plus longtemps
qu'elle n'est développée, et qu'une part significative du travail d'une équipe
consiste à maintenir en condition opérationnelle ce qui existe déjà. Cela a
modifié ma manière de concevoir~: je me demande désormais, en écrivant une
fonctionnalité, comment elle sera diagnostiquée le jour où elle
dysfonctionnera.

\section{Lien avec la formation ENSIM}
\label{sec:lien-ensim}

Plusieurs enseignements suivis à l'ENSIM se sont révélés directement
mobilisables, mais rarement sous la forme dans laquelle je les avais appris.

Le cours de bases de données relationnelles est celui dont l'apport a été le
plus direct. La modélisation et les formes normales m'ont fourni le cadre
d'analyse qui m'a permis de qualifier le problème du modèle existant, et de
proposer une solution argumentée plutôt qu'une intuition. Ce que le stage a
ajouté, c'est la dimension temporelle et opérationnelle que l'exercice
académique ignore~: normaliser une base vide est un problème de conception,
normaliser une base en production est un problème de migration, qui suppose des
scripts réversibles, des contrôles préalables et une reprise des données
existantes.

Le génie logiciel, la conception objet et le développement web ont constitué le
socle de mon travail sur l'architecture en couches de l'application. La
séparation des responsabilités, la notion de contrat entre composants et les
patrons de conception vus en cours m'ont permis de me repérer rapidement dans
une base de code que je n'avais pas écrite, et l'architecture n-tiers étudiée en
formation s'est concrétisée dans le découpage réel de la solution, en projets
distincts dont chacun a une responsabilité assignée. Je n'avais pratiqué ni
Angular ni~.NET avant le stage~; c'est la compréhension de ces concepts
sous-jacents, plutôt que la connaissance d'un framework particulier, qui m'a
permis de monter en compétence sans blocage majeur, ce qui confirme l'intérêt
d'une formation qui enseigne des principes plutôt que des outils. Le stage m'a
appris en revanche ce qui s'enseigne difficilement en cours~: qu'une
architecture existante porte des compromis, et qu'un développeur qui la reprend
doit d'abord comprendre pourquoi ces compromis ont été faits avant de les
corriger.

Ma spécialité, Interaction Personnes Systèmes, a trouvé une application moins
attendue mais réelle. La dette ergonomique identifiée dans l'application
existante, l'absence de retour visuel pendant les chargements, la difficulté à
retrouver une information dans de longues listes, relève exactement des
questions traitées en formation. Le travail sur la lisibilité de l'historique
des actions en est l'exemple le plus net~: le problème n'était pas de tracer les
modifications, mais de les présenter à une granularité compréhensible pour
l'utilisateur plutôt qu'à celle du système. Les exigences d'accessibilité
remontées par les analyses de qualité ont, de la même manière, prolongé une
sensibilité acquise en cours.

Les enseignements de gestion de projet et de méthodes agiles m'ont permis
d'arriver en connaissant le vocabulaire et la logique des rituels. Leur pratique
quotidienne m'a toutefois appris quelque chose que le cours ne transmet pas~:
l'agilité en entreprise n'est pas un cadre appliqué à la lettre, mais une
organisation ajustée à la réalité d'une équipe, avec des rituels dont chacun
a une fonction précise.

Trois domaines ont en revanche mis en évidence un écart réel entre ma formation
et le terrain. Le cloud, d'abord~: si les notions de réseaux et de systèmes
m'ont aidé à comprendre les mécanismes sous-jacents, je n'avais aucune pratique
d'une plateforme comme Azure, et j'ai dû acquérir en autonomie la logique des
services managés, de l'authentification par certificat et de la séparation des
environnements. La mise en production, ensuite~: c'est probablement l'écart le
plus important, car aucun projet académique ne m'avait confronté à l'idée qu'une
erreur ne se corrige pas par une nouvelle exécution du programme, mais engage
des données et des utilisateurs réels. La montée de version d'un socle
technique, enfin, est une opération que je n'avais jamais eu l'occasion de
pratiquer, et dont j'ignorais qu'elle constitue une activité à part entière du
métier.

Ce stage m'a également révélé deux compétences qu'il me reste à construire. La
première concerne les tests automatisés~: l'application dont j'avais la charge
n'en disposait pas côté serveur, et j'ai validé mon travail par des tests
manuels et par la recette fonctionnelle. Cette situation m'a fait mesurer, par
son absence, ce qu'une couverture de tests apporte en confiance lors d'une
refonte, mais elle ne m'a pas permis d'en développer la pratique. La seconde
concerne le déploiement automatisé~: j'ai travaillé avec une chaîne
d'intégration continue qui compile et analyse le code, sans intervenir sur son
paramétrage, et les déploiements auxquels j'ai participé restaient largement
manuels. Ce sont les deux axes sur lesquels je souhaite progresser en priorité.

Au terme de cette expérience, ma représentation du métier d'ingénieur
développeur a changé. Je l'envisageais essentiellement comme une activité de
production de code~; je le comprends aujourd'hui comme un travail d'analyse et
d'arbitrage, où la difficulté réside moins dans l'écriture d'une solution que
dans l'identification du bon problème et dans la capacité à le formuler
suffisamment clairement pour qu'une équipe décide de le traiter.
